% ============================================
% Преамбула отчёта НИР
% Проект: sdCloud — AI-ассистент для генерации
% системно-динамических моделей
% ============================================

\documentclass[14pt,a4paper]{extreport}

% ---------- Кодировка и шрифты ----------
\usepackage[utf8]{inputenc}
\usepackage[T2A]{fontenc}

% ---------- Язык ----------
\usepackage[russian]{babel}

% ---------- Геометрия страницы (ГОСТ) ----------
\usepackage{geometry}
\geometry{
    a4paper,
    left=30mm,
    right=15mm,
    top=20mm,
    bottom=20mm,
    footskip=15mm,
    headsep=10mm
}

% ---------- Интервалы ----------
\usepackage{setspace}
\onehalfspacing

% ---------- Абзацный отступ ----------
\usepackage{indentfirst}
\setlength{\parindent}{1.25cm}

% ---------- Математика ----------
\usepackage{amsmath,amssymb,amsthm}

% ---------- Графика ----------
\usepackage{graphicx}
\usepackage{tikz}
\usetikzlibrary{shapes.geometric, arrows.meta, positioning, calc, fit, backgrounds, decorations.pathreplacing}

% ---------- Таблицы ----------
\usepackage{booktabs,longtable,multirow,array,colortbl}
\usepackage{tabularx}
\newcolumntype{L}[1]{>{\raggedright\let\newline\\\arraybackslash\hspace{0pt}}m{#1}}
\newcolumntype{C}[1]{>{\centering\let\newline\\\arraybackslash\hspace{0pt}}m{#1}}
\newcolumntype{R}[1]{>{\raggedleft\let\newline\\\arraybackslash\hspace{0pt}}m{#1}}

% ---------- Листинги кода ----------
\usepackage{listings}
\usepackage{xcolor}

% ---------- Гиперссылки ----------
\usepackage[
    colorlinks=true,
    linkcolor=black,
    citecolor=black,
    urlcolor=blue,
    bookmarks=true,
    bookmarksopen=true,
    pdftitle={sdCloud: AI-ассистент для генерации системно-динамических моделей},
    pdfauthor={ИТМО}
]{hyperref}

% ---------- Колонтитулы ----------
\usepackage{fancyhdr}
\pagestyle{fancy}
\fancyhf{}
\fancyfoot[C]{\thepage}
\renewcommand{\headrulewidth}{0pt}
\renewcommand{\footrulewidth}{0pt}
\fancypagestyle{plain}{
    \fancyhf{}
    \fancyfoot[C]{\thepage}
    \renewcommand{\headrulewidth}{0pt}
    \renewcommand{\footrulewidth}{0pt}
}

% ---------- Заголовки ----------
\usepackage{titlesec}
\titleformat{\chapter}[display]
    {\normalfont\Large\bfseries\centering}
    {\chaptertitlename\ \thechapter}{14pt}{\Large}
\titlespacing*{\chapter}{0pt}{-20pt}{20pt}

\titleformat{\section}
    {\normalfont\large\bfseries}
    {\thesection}{1em}{}
\titleformat{\subsection}
    {\normalfont\normalsize\bfseries}
    {\thesubsection}{1em}{}
\titleformat{\subsubsection}
    {\normalfont\normalsize\itshape}
    {\thesubsubsection}{1em}{}

% ---------- Нумерация ----------
\renewcommand{\thechapter}{\arabic{chapter}}
\renewcommand{\thesection}{\thechapter.\arabic{section}}
\renewcommand{\thesubsection}{\thesection.\arabic{subsection}}
\renewcommand{\thesubsubsection}{\thesubsection.\arabic{subsubsection}}

% ---------- Названия элементов ----------
\addto\captionsrussian{
    \renewcommand{\contentsname}{СОДЕРЖАНИЕ}
    \renewcommand{\chaptername}{ГЛАВА}
    \renewcommand{\figurename}{Рис.}
    \renewcommand{\tablename}{Таблица}
    \renewcommand{\bibname}{СПИСОК ЛИТЕРАТУРЫ}
    \renewcommand{\lstlistingname}{Листинг}
}

% ---------- Подписи к рисункам и таблицам ----------
\usepackage{caption}
\captionsetup[figure]{labelsep=endash, justification=centering, font=small}
\captionsetup[table]{labelsep=endash, justification=justified, font=small, singlelinecheck=false}
\captionsetup[lstlisting]{labelsep=endash, font=small}

% ---------- Стили листингов ----------
\lstset{
    basicstyle=\small\ttfamily,
    keywordstyle=\color{blue}\bfseries,
    commentstyle=\color{gray}\itshape,
    stringstyle=\color{red},
    numbers=left,
    numberstyle=\tiny\color{gray},
    numbersep=8pt,
    frame=single,
    breaklines=true,
    breakatwhitespace=true,
    tabsize=4,
    showstringspaces=false,
    inputencoding=utf8,
    extendedchars=true,
    literate=
        {а}{\cyra}1 {б}{\cyrb}1 {в}{\cyrv}1 {г}{\cyrg}1 {д}{\cyrd}1
        {е}{\cyre}1 {ё}{\cyryo}1 {ж}{\cyrzh}1 {з}{\cyrz}1 {и}{\cyri}1
        {й}{\cyrishrt}1 {к}{\cyrk}1 {л}{\cyrl}1 {м}{\cyrm}1 {н}{\cyrn}1
        {о}{\cyro}1 {п}{\cyrp}1 {р}{\cyrr}1 {с}{\cyrs}1 {т}{\cyrt}1
        {у}{\cyru}1 {ф}{\cyrf}1 {х}{\cyrh}1 {ц}{\cyrc}1 {ч}{\cyrch}1
        {ш}{\cyrsh}1 {щ}{\cyrshch}1 {ъ}{\cyrhrdsn}1 {ы}{\cyry}1 {ь}{\cyrsoftsign}1
        {э}{\cyrie}1 {ю}{\cyryu}1 {я}{\cyrya}1
        {А}{\CYRA}1 {Б}{\CYRB}1 {В}{\CYRV}1 {Г}{\CYRG}1 {Д}{\CYRD}1
        {Е}{\CYRE}1 {Ё}{\CYRYO}1 {Ж}{\CYRZH}1 {З}{\CYRZ}1 {И}{\CYRI}1
        {Й}{\CYRISHRT}1 {К}{\CYRK}1 {Л}{\CYRL}1 {М}{\CYRM}1 {Н}{\CYRN}1
        {О}{\CYRO}1 {П}{\CYRP}1 {Р}{\CYRR}1 {С}{\CYRS}1 {Т}{\CYRT}1
        {У}{\CYRU}1 {Ф}{\CYRF}1 {Х}{\CYRH}1 {Ц}{\CYRC}1 {Ч}{\CYRCH}1
        {Ш}{\CYRSH}1 {Щ}{\CYRSHCH}1 {Ъ}{\CYRHRDSN}1 {Ы}{\CYRY}1 {Ь}{\CYRSOFTSIGN}1
        {Э}{\CYRIE}1 {Ю}{\CYRYU}1 {Я}{\CYRYA}1
}

% ---------- Окружения ----------
\newtheorem{definition}{Определение}[chapter]
\newtheorem{theorem}{Теорема}[chapter]

% ---------- Пакеты для алгоритмов ----------
\usepackage{algorithm}
\usepackage{algpseudocode}
\floatname{algorithm}{Алгоритм}

% ---------- Дополнительные пакеты ----------
\usepackage{enumitem}
\setlist[itemize]{leftmargin=1.5cm,itemsep=2pt}
\setlist[enumerate]{leftmargin=1.5cm,itemsep=2pt}
\usepackage{float}
\usepackage{csquotes}
\usepackage{url}
\usepackage{breakurl}
\usepackage{footnote}
\usepackage{microtype}
\usepackage{etoolbox}

% ---------- Счётчики ----------
\setcounter{secnumdepth}{3}
\setcounter{tocdepth}{2}

% ---------- Компактные списки ----------
\let\olditemize\itemize
\let\oldenumerate\enumerate
\renewenvironment{itemize}{
    \olditemize[topsep=4pt,parsep=2pt]
}{\endolditemize}
\renewenvironment{enumerate}{
    \oldenumerate[topsep=4pt,parsep=2pt]
}{\endoldenumerate}
